% Source: tex/chapters/ch05/10_新興技術のテストと新興技術によるテスト.tex
\subsubsection{新興技術を使ってテストする}
機械学習は、テストケース設計、オラクル問題、テスト評価、優先順位付け、ミューテーションテスト自動化、欠陥分類、トリアージに使われる。大量データから傾向を抽出し、保守労力を減らす可能性がある。

ブロックチェーンは、分散した関係者が信頼できる形でテスト活動を調整する手段になり得る。改ざん耐性、監査可能性、分散データ管理、要求順守の自動確認に役立つ可能性がある。クラウドは、大規模シミュレーションや弾力的資源を必要とするテスト環境として利用できる。

シミュレーションは、危険な状況、災害、復旧、特定振る舞いを評価するための重要な手段である。自動車や組込み領域では、実信号を SUT に与えながら現実を模擬するハードウェアインザループテストが使われる。クラウドソーシングテストは、多様な利用者、場所、端末、条件での検出力を得るために使われるが、社内検証の代替ではなく補完である。
